\uuid{lIdq}
\titre{Matrice d'une application linéaire}
\auteur{hanani ; mzari}
\organisation{exo7}
\licence{CC BY-NC-SA 4.0}
\datecreate{2020-01-20}
\niveau{L1}
\module{Algèbre}
\chapitre{Matrice}
\sousChapitre{Matrice et application linéaire}
\difficulte{2}
\theme{Matrice d'une application linéaire}
\exo7qcmid{lille-2.matrice-application.010}

\contenu{
\begin{question}
\qclassification{108.03}
\qtags{motcle=Matrice d'une application linéaire}
Soit $u_1=(1,0,0), \;  u_2=(1,1,0), \;  u_3=(0,1,1), \; v_1=(1,1), \; v_2=(1,-1)$  et $f$ l'application linéaire définie par :
$$\begin{array}{rccc}f:&\Rr^3 &\to& \Rr^2\\
& (x,y,z)&\to &(x+y,x-z).  \end{array}$$
Quelles sont les assertions vraies ?
\begin{answers}
\good{$\{ u_1, u_2, u_3\}$ est une base de $\Rr^3$.}
\good{$\{ v_1, v_2\}$ est une base de $\Rr^2$.}
\bad{La matrice de $f$ par rapport aux bases  $\{ u_1, u_2, u_3\}$ et $\{ v_1, v_2\}$ est :
$\displaystyle \frac{1}{2}\left(\begin{array}{rc}
2&0\\ 3&1\\ 0&2\end{array}\right)$.}
\good{La matrice de $f$ par rapport aux bases  $\{ u_1, u_2, u_3\}$ et $\{ v_1, v_2\}$ est :
$$\displaystyle
\frac{1}{2} \left(\begin{array}{rcc}
2&3&0\\ 0&1&2\end{array}\right).$$}
\end{answers}
\begin{explanations} On vérifie que $\{ u_1, u_2, u_3\}$ est une base de $\Rr^3$ et que $\{ v_1, v_2\}$ est une base de $\Rr^2$. La matrice de $f$  par rapport à ces  bases    est la matrice  dont la $j$ième colonne est constituée des coordonnées de  $f(u_j)$ dans la base $\{ v_1, v_2\}$. Cette matrice est : $\displaystyle \frac{1}{2} \left(\begin{array}{rcc} 2&3&0\\ 0&1&2\end{array}\right)$.
\end{explanations}
\end{question}
}
